\documentclass[preprint,12pt]{elsarticle}

% Common packages
\usepackage{graphicx}
\usepackage{amssymb}
\usepackage{amsmath}
\usepackage{hyperref}
\usepackage{subcaption}
\usepackage{booktabs}
\usepackage{multirow}
\usepackage{tikz}
\usepackage{pgfplots}
\usepgfplotslibrary{groupplots}
\pgfplotsset{compat=1.18}

% Optional packages
% \usepackage{amsthm}
% \usepackage{lineno}

\journal{Journal Name}

\begin{document}

\begin{frontmatter}

\title{ADEPT: Query-Conditioned Event Graph Forecasting with Adaptive Delta-Kernel Experts for \\ Irregular Time Series}

\fntext[fn1]{Nguyen Quoc Anh and Trung Nam Nguyen contributed equally to this work and share co-first authorship.}

\cortext[cor1]{Corresponding author.}

\author[label1]{Trung Nam Nguyen\fnref{fn1}}
\ead{trungnn@vaa.edu.vn}

\author[label2]{Nguyen Quoc Anh\fnref{fn1}}
\ead{anh.quocnguyen@hitachids.com}

\author[label3]{Ha Xuan Son}
\ead{ha.son@rmit.edu.vn}

\author[label4,label5]{Nguyen~Ngoc~Phien\corref{cor1}}
\ead{nguyenngocphien@tdtu.edu.vn}

\author[label6]{Trung Phan-Hoang Tuan}
\ead{trungpht@fe.edu.vn}

\author[label6]{Nguyen Tien Thuan}
\ead{ThuanNTCE181024@fpt.edu.vn}

\author[label6]{Nguyen Gia Chan}
\ead{channgce181288@fpt.edu.vn}

\author[label6]{Le Huu Khoa}
\ead{khoalhce181099@fpt.edu.vn}

\affiliation[label1]{
  organization={Faculty of Business Administration, Vietnam Aviation Academy},
  city={Ho Chi Minh City},
  country={Vietnam}
}

\affiliation[label2]{
  organization={Department of Product Engineering and Embedded Systems, A\&I CoE, Hitachi Digital Services},
  city={Ho Chi Minh City},
  state={District 12},
  country={Vietnam}
}

\affiliation[label3]{
  organization={Department of Blockchain-Enabled Business, RMIT University},
  city={Ho Chi Minh City},
  state={District 7},
  country={Vietnam}
}

\affiliation[label4]{
  organization={Center for Applied Information Technology, Ton Duc Thang University},
  city={Ho Chi Minh City},
  country={Vietnam}
}

\affiliation[label5]{
  organization={Faculty of Information Technology, Ton Duc Thang University},
  city={Ho Chi Minh City},
  country={Vietnam}
}

\affiliation[label6]{
  organization={FPT University, Can Tho Campus},
  city={Can Tho},
  country={Vietnam}
}

\begin{abstract}
Irregular multivariate time-series forecasting is fundamentally a retrieval
problem: each future value asks which asynchronous historical events, variables,
and time gaps should be trusted. Dense imputation and patching make this setting
easier to process, but they can hide the event structure that determines the
answer. We introduce ADEPT, a query-conditioned event graph forecaster with
adaptive delta-kernel experts. ADEPT encodes observations as value-time-variable
event tokens, contextualizes them with transformer memory, learns cross-variable
relations, and retrieves evidence for each future query through continuous-time
delta kernels. Its distinguishing component is an end-to-end expert decoder
that gates between the neural query forecast and history-derived anchors for
same-variable persistence, exponential memory, local trend, graph-aware context,
and global entity state, without external model predictions, test labels, or
post hoc routing. On four standard IMTS benchmarks, ADEPT delivers the best
standard-window MAE in every dataset, improving over the strongest baseline by
2.6--22.2\% and leading MSE on three datasets. Across changing lookback and
horizon settings, the advantage remains absolute-error driven rather than
window-specific. Through perturbation and ablation evidence, the finding is
that accuracy depends on using real values, temporal placement, event memory,
and adaptive experts, in whichs ADEPT gains by retrieving the right events, not by
exploiting masks or persistence.
\end{abstract}

\begin{keyword}
Irregular time series \sep Event graph neural networks \sep Query-conditioned forecasting \sep Adaptive delta-kernel experts
\end{keyword}

\end{frontmatter}

% \linenumbers

\newcommand{\method}{ADEPT}
\newcommand{\best}[1]{\textbf{#1}}
\newcommand{\second}[1]{\underline{#1}}

\section{Introduction}
\label{sec:introduction}

Irregular multivariate time series (IMTS) arise when variables are observed at
different times, under different sampling rates, and often for reasons that are
themselves informative. Clinical records follow care decisions, climate
stations report sparse seasonal measurements, and wearable sensors produce
event-like trajectories rather than clean grids. Forecasting in this setting is
therefore not simply a sequence-to-vector problem. Each future target asks a
more specific question: which past events, variables, and time gaps should be
trusted for this particular future time--variable query?

Existing IMTS methods attack different sides of this question. Missingness-aware
recurrence, interpolation, time attention, and set encoders make irregular
observations easier to process \citep{grud,shukla2019interp,mtan,seft}.
Continuous-time models represent latent dynamics between observations
\citep{neuralode2,latentODE,neuralCDE,cru}. Graph and patch models preserve or
regularize cross-variable structure under sparse observation patterns
\citep{raindrop2,gneuralflow,grafiti,tPatchGNN,apn}. These advances are
substantial, but they often shift the forecasting decision into an imputed
grid, a latent trajectory, a patch representation, or a mostly neural final
head. The final query still has to recover a simple but difficult balance:
same-variable persistence may be decisive for one target, graph-aware evidence
for another, and a local trend or entity-level context for a third.

We introduce \method{}, a query-conditioned event graph forecaster designed
around that final evidence-selection step. Observed measurements remain event
tokens with value, variable identity, and continuous time information. A
transformer memory contextualizes the event set, an adaptive variable graph
learns cross-channel relations, and each future query retrieves evidence through
both cross-attention and learned continuous-time delta kernels. The key
difference is the adaptive delta-kernel expert decoder: instead of asking a
single neural head to rediscover all local numerical structure, \method{} gates
between the neural query forecast and differentiable history-derived experts for
same-variable latest value, exponential memory, local trend, graph-aware memory,
and global entity state. A bounded residual then lets the neural decoder correct
the gated forecast.

This design keeps the protocol honest. The experts are not external baseline
predictions, not a hand-tuned ensemble, and not a test-label calibration
mechanism. They are computed only from the observed history of the current
example, their weights are produced by the model, and the whole system is
trained end-to-end with masked forecasting losses. The goal is to preserve the
strong numerical anchors that classical forecasting teaches us to respect while
retaining the flexibility of neural event memory.

Under the evaluated IMTS protocol, this combination changes the empirical
profile of the model. \method{} achieves the best standard-window MAE on P12,
MIMIC, USHCN, and HumanActivity, improving over the strongest baseline by
2.6--22.2\%, and it obtains the best standard-window MSE on three of the four
datasets. Its advantage is not confined to one canonical split: across the
lookback--horizon sweep, \method{} remains consistently strongest on MAE.
Perturbations and ablations further show that the gains require real observed
values, correct timing, event-memory retrieval, and the adaptive expert gate.
The contributions are:
\begin{itemize}
    \item We formulate IMTS forecasting as query-conditioned event retrieval
    and propose \method{}, combining transformer event memory, adaptive
    variable graphs, delta-kernel retrieval, and an end-to-end adaptive expert
    decoder.
    \item We show state-of-the-art standard-window MAE on all four benchmarks,
    best MSE on three, and a 2.6--22.2\% MAE reduction over the strongest
    baseline under matched masks, scaling, windows, and seeds.
    \item We diagnose the mechanism through window sweeps, temporal
    perturbations, entity-level uncertainty, and ablations, showing robust
    absolute-error gains while identifying the remaining squared-error boundary
    on HumanActivity and selected USHCN windows.
\end{itemize}

\section{Related Work}
\label{sec:related-work}

\paragraph{Regular-grid forecasting and local anchors}
Classical forecasting remains relevant to IMTS because it clarifies what a
strong forecast must preserve: a useful local level, trend, seasonality, and
uncertainty structure. ARIMA and state-space models provide transparent
temporal dynamics \citep{BoxJenkins1970,Harvey1989,Hamilton1994,DurbinKoopman2012}.
Exponential smoothing families give strong local anchors for level and trend
\citep{Holt1957,Winters1960,Gardner1985,Hyndman2002,Hyndman2008}. Seasonal
decomposition and forecasting competitions reinforce the same lesson: simple
models can remain difficult to beat when they encode the right inductive bias
\citep{Cleveland1990STL,Assimakopoulos2000Theta,Makridakis2018M4,Makridakis2022M5}.
Machine-learning regressors brought nonlinear residual modeling and strong
tabular learners into forecasting pipelines \citep{Drucker1997SVR,Smola2004SVR,Breiman2001RandomForest,Friedman2001GBM}.
Boosted trees and hybrid statistical--neural systems further show that learned
corrections work best when they are attached to stable baseline structure
\citep{Chen2016XGBoost,Ke2017LightGBM,Zhang2003HybridARIMAANN,KhasheiBijari2011Hybrid,MonteroManso2020FFORMA,Smyl2020ESRNN}.
These models do not solve asynchronous multivariate forecasting, but they
motivate one design principle behind \method{}: a neural IMTS model should not
throw away same-variable persistence, trend, and graph-aware local evidence
while learning a richer representation. Deep forecasting moved from recurrent sequence models to attention-based
sequence transduction \citep{Hochreiter1997LSTM,Cho2014GRU,Sutskever2014Seq2Seq,Bahdanau2015Attention}.
Transformer, probabilistic, basis-expansion, and interpretable multi-horizon
models broadened that line \citep{Vaswani2017Transformer,Salinas2020DeepAR,Oreshkin2020NBEATS,Lim2021TFT}.
Long-horizon architectures then improved regular-grid forecasting through
sparse attention, decomposition, and frequency-domain structure
\citep{informer,Wu2021Autoformer,Zhou2022FEDformer,Wu2023TimesNet}. Other
models emphasized patching, linear or MLP mixing, and cross-variable
tokenization \citep{Nie2023PatchTST,Zeng2023LTSFLinear,Das2023TiDE,Ekambaram2023TSMixer}.
Graph and hypergraph variants made cross-channel structure more explicit
\citep{Crossformer2023,FourierGNN2023,AdaMSHyper2024,iTransformer2024,TimeMixer2024}.
Most recently, state-space and foundation-model approaches have made
large-scale pretraining and efficient long-context sequence modeling central to
time-series forecasting \citep{Mamba2024,MambaTS2024,Moirai2024,TimesFM2024,Chronos2024,LagLlama2023}.
This line is powerful, but its interface is usually an aligned tensor with a
fixed input window and output horizon. When such models are applied to IMTS,
irregularity is often handled before modeling through resampling, imputation,
mask channels, or patches. That preprocessing can be useful, but it changes the
main question from ``which historical event supports this query?'' to ``what
regular representation should we forecast from?''

\paragraph{Irregular representations and continuous time}
Early neural IMTS models addressed the irregular sampling problem by adding
time-gap and missingness information directly to sequence or interpolation
mechanisms. GRU-D uses masking and elapsed-time decay to make recurrence aware
of missingness patterns \citep{grud}. Interpolation-prediction networks and
mTAN instead build explicit interpolation or time-aware attention layers before
prediction \citep{shukla2019interp,mtan}. Set-based modeling offers another
route: SeFT treats observations as unordered event sets, reducing dependence on
an arbitrary dense grid \citep{seft}. These approaches were important because
they made irregularity part of the model, yet they still tend to summarize the
history before the final forecast is queried. Continuous-time latent dynamics provide a more principled way to model
irregular observation times. Neural ODE and Latent ODE models evolve hidden
states in continuous time, while Neural CDEs represent the observed path as a
control signal \citep{neuralode2,latentODE,neuralCDE}. GRU-ODE-Bayes and CRU
combine continuous evolution with update mechanisms suited to partially
observed clinical time series \citep{DeBrouwer2019GRUODE,cru}. Later variants
added attention, learned interpolation/extrapolation, flows, or transformer
relations in continuous time \citep{ANCDE2021,EXIT2022,neuralflows2,ContiFormer2023}.
Diffusion-based imputation also improved the quality of missing-value
reconstruction and can support forecasting through an impute-then-predict
pipeline \citep{CSDI2021}. The strength of this family is temporal coherence;
its risk is over-compression. A latent trajectory or imputed grid may be smooth
and useful, but the downstream forecast can lose direct access to the particular
events whose values, recency, and variables matter most for a future target.

\paragraph{Graphs, patches, and event-query forecasting}
Recent IMTS forecasting work moves closer to the raw observation protocol by
representing interactions among variables, observations, and forecast targets.
Raindrop and GNeuralFlow use graph-guided propagation to capture cross-channel
relationships under irregular observation patterns \citep{raindrop2,gneuralflow}.
PrimeNet studies pretraining for irregular multivariate series, while
Warpformer introduces multi-scale warping for irregular clinical signals
\citep{primenet,warpformer}. GraFITi is especially close to an event-query view:
it converts observed and target sparsity into a graph and casts forecasting as
edge prediction \citep{grafiti}. These models preserve more of the original
sampling structure than dense-grid forecasters, but the final numerical decoder
still has to decide how much to trust same-variable recency, cross-variable
evidence, and learned latent context. Patch-based IMTS models provide an alternative answer. tPatchGNN showed that
patching irregular histories can make graph forecasting efficient and strong,
and APN further improved this direction through adaptive patching
\citep{tPatchGNN,apn}. HyperIMTS, Hi-Patch, and ASTGI push graph construction
further by using hyperedges, hierarchical patch graphs, or adaptive
spatio-temporal neighborhoods \citep{HyperIMTS2025,HiPatch2025,ASTGI2026}.
DNBNet revisits another important angle: basis-function representations can
suffer from timestamp-density bias and limited adaptability, so irregular
forecasting also depends on how histories are numerically summarized
\citep{DNBNet2026}. Together, these models show that structure matters, but
they also expose a tension. Patches and basis responses improve efficiency and
regularity, while raw event graphs preserve fidelity; neither alone guarantees
that the final forecast will keep the strongest local numerical anchors.

\paragraph{Positioning of our \method{}}
\method{} is designed for this gap between representation and final numerical
forecasting. Unlike dense-grid and foundation forecasters, it keeps the input as
events and decodes explicit future time--variable queries. Unlike purely
continuous-time latent models, it leaves the decoder direct access to retrieved
observations through cross-attention and learned delta kernels. Unlike graph or
patch models whose final head is mostly neural, it exposes the query decoder to
history-derived same-variable, trend, graph-aware, and global experts. The
experts are not external baselines and are not selected using test labels; they
are differentiable functions of the observed history and are gated end-to-end.
The contribution is therefore not just a new encoder or a new graph
construction. It is a protocol-preserving event-query architecture in which
adaptive neural memory and explicit history anchors cooperate at the point
where IMTS forecasting is hardest: deciding what evidence each future
time--variable query should trust.

\section{Methodology}
\label{sec:method}

\subsection{Architectural Proposal}
\label{sec:architecture}

\begin{figure}[htbp]
\centering
\makebox[\textwidth]{%
  \includegraphics[width=1.5\textwidth]{ADEPT (Processing)/architecture.png}}
\caption{High-level design of the proposed \method{} architecture. The
implementation uses event-token encoding, transformer memory, adaptive variable
graph mixing, continuous-time delta-kernel retrieval, and an adaptive expert
query decoder. Fixed training settings are reported in
Appendix~\ref{app:adept-config}.}
\label{architecture}
\end{figure}

Notation used in this subsection is summarized in
Table~\ref{tab:adept-symbols} of Appendix~\ref{app:adept-symbols}. For each entity, \method{} preserves the
observed history as an event set and decodes the masked future targets as
explicit time--variable queries:
\begin{equation}
    \mathcal E_b = \{(t_i,v_i,\widetilde x_i)\}_{i=1}^{N_b},
\end{equation}
\begin{equation}
    \mathcal Q_b=\{(\tau_j,u_j)\}_{j=1}^{M_b}.
\end{equation}
Losses and metrics are computed only on active target-mask entries. This
formulation keeps the modeling target aligned with the benchmark protocol:
forecasts are requested values at specific future times and variables, not a
dense vector produced after resampling.

\paragraph{Event tokenization}
Each finite observation is mapped to a neural token by adding a variable
embedding, a value projection, and a continuous-time Fourier projection:
\begin{equation}
    h_i^{(0)} = e_{v_i} + f_x(\widetilde x_i) + f_t(t_i),
\end{equation}
so the model sees time, variable identity, and measured value before any
regular-grid conversion. A transformer encoder then contextualizes the padded
event memory:
\begin{equation}
    H_b = \operatorname{TransformerEncoder}\bigl(h_1^{(0)},\ldots,h_{N_b}^{(0)}\bigr).
\end{equation}

\paragraph{Adaptive variable graph}
Cross-variable structure is learned from the same variable embeddings used by
events and queries:
\begin{equation}
    A_{u,v} = \operatorname{softmax}_v\left(\frac{\bar e_u^\top \bar e_v}{\sqrt d}\right).
\end{equation}
Each row of $A$ is a query-conditioned distribution over source variables. The
graph context $a_u=\sum_v A_{u,v}e_v$ is projected and used in query
construction, temporal retrieval, and the graph-aware expert. This gives the
model dataset-specific channel relations without requiring a fixed clinical,
climate, or sensor graph.

\paragraph{Query-conditioned temporal retrieval}
For each future query $(\tau_j,u_j)$, \method{} builds a query token from the
requested variable embedding, Fourier time features, and graph context. The
query attends to transformer event memory using multi-head cross-attention. In
parallel, $K$ learned continuous-time delta kernels retrieve recency-sensitive
memory:
\begin{equation}
    r_{j,k}=\sum_i \alpha_{j,i,k}H_{b,i},\qquad
    \alpha_{j,i,k}=\operatorname{softmax}_i\left(
    -\frac{|\tau_j-t_i|}{s_k}
    + \beta\,\mathbb{1}[u_j=v_i]
    + \gamma\log A_{u_j,v_i}\right).
\end{equation}
The retrieval score combines temporal proximity, same-variable bias, and
adaptive graph affinity. The decoder therefore receives both global attention
context and graph-aware temporal memory at multiple learned time scales.

\paragraph{Adaptive delta-kernel expert decoder}
The adaptive expert decoder is designed to make the neural query state
numerically precise without leaving the end-to-end training protocol. First, a
deep query state $d_j$ is formed from the query token, cross-attention context,
delta-kernel context, graph context, and a lookback-normalized query gap through
a feed-forward decoder backbone. A neural head produces a base forecast
$z^{\mathrm{nn}}_j$. In addition, the decoder constructs five scalar experts from the observed
history:
\begin{align}
z^{\mathrm{last}}_j &= \text{soft latest-value estimate for variable }u_j,\\
z^{\mathrm{ema}}_j &= \sum_{i:v_i=u_j} \omega_{j,i}^{u}\widetilde x_i,\\
z^{\mathrm{trend}}_j &= z^{\mathrm{last}}_j +
    \widehat m_j(\tau_j-t^{\mathrm{last}}_{u_j}),\\
z^{\mathrm{graph}}_j &= \sum_i \omega_{j,i}^{g}\widetilde x_i,\\
z^{\mathrm{global}}_j &= \sum_i \omega_{j,i}^{0}\widetilde x_i.
\end{align}
The same-variable and global experts use normalized recency weights, the
graph-aware expert uses a recency--graph kernel, and the trend expert
extrapolates a clipped local slope when enough same-variable evidence exists.
When same-variable history is unavailable, the same-variable experts fall back
to graph-aware memory and the gate receives an availability indicator. The final forecast is a gated mixture with a bounded residual:
\begin{equation}
    \pi_j=\operatorname{softmax}\left(W_\pi[d_j,z_j,m_j,g_j]\right),
\end{equation}
\begin{equation}
    \widehat y_j
    = \sum_{\ell}\pi_{j,\ell}z_{j,\ell}
    + 0.05\tanh\bigl(r_\theta(d_j,z_j,m_j,g_j)\bigr).
\end{equation}
The gate sees the deep query state, expert values, same-variable availability,
and query recency gap. History-derived expert values are clipped in scaled
space before gating. This structure gives the model access to strong local
persistence and graph-aware retrieval while letting the neural residual handle
bounded nonlinear corrections.

% \paragraph{Training}
% The model is optimized on the training split using a masked objective combining
% MSE, smooth-$L_1$, and MAE terms:
% \begin{equation}
%     \mathcal L = 0.50\mathcal L_{\mathrm{MSE}}
%     +0.20\mathcal L_{\mathrm{Smooth}L_1}
%     +0.30\mathcal L_{\mathrm{MAE}}.
% \end{equation}
% The validation score used for early stopping is
% $\mathrm{MSE}+0.25\,\mathrm{MAE}$. This places the main selection pressure on
% squared error while keeping the decoder aligned with the benchmark MAE target.
% All reported \method{} benchmark results are averaged over seeds 2024, 2025, and
% 2026.

\subsection{Benchmark Protocol}
\label{sec:protocol}

\paragraph{Datasets}
We evaluate on four widely used IMTS forecasting benchmarks: PhysioNet Challenge
2012 (P12), MIMIC-III, USHCN, and HumanActivity. The datasets differ in
sampling density, number of variables, temporal scale, and target availability.
P12 and MIMIC contain intensive-care measurements, USHCN contains irregular
climate station observations, and HumanActivity contains wearable-sensor
trajectories.

\paragraph{Baselines}
We compare \method{} with representative neural IMTS models under the same
benchmark protocol. The primary neural comparators are APN \citep{apn},
GraFITi \citep{grafiti}, and tPatchGNN \citep{tPatchGNN}, which are the
strongest recent reproduced baselines under this protocol. The full table also
reports earlier neural baselines from the reference protocol, including
PrimeNet \citep{primenet}, NeuralFlows \citep{neuralflows2}, CRU \citep{cru},
mTAN \citep{mtan}, SeFT \citep{seft}, GNeuralFlow \citep{gneuralflow}, GRU-D
\citep{grud}, and Raindrop \citep{raindrop2}. Neural baselines are evaluated
using their released configurations with validation-based hyperparameter tuning,
and performance is averaged across seeds when available.

\paragraph{Training and Metrics}
All methods are compared under the same dataset splits, lookback--horizon
windows, input masks, target masks, and benchmark scaling. Baselines use their
released training objectives, optimizers, and validation-based model selection,
so each method is evaluated under its intended configuration rather than being
forced into a shared optimizer or loss. \method{} uses one fixed training
recipe across the benchmark runs; its loss weights are not tuned separately for
individual datasets, windows, seeds, or test outcomes. Training uses only the
training split with a masked objective combining MSE, smooth-$L_1$, and MAE:
\begin{equation}
    \mathcal L = 0.50\mathcal L_{\mathrm{MSE}}
    +0.20\mathcal L_{\mathrm{Smooth}L_1}
    +0.30\mathcal L_{\mathrm{MAE}}.
\end{equation}
Checkpoint selection uses validation data only, with
$\mathrm{MSE}+0.25\,\mathrm{MAE}$ as the pre-specified \method{} validation
score. No method uses test labels for hyperparameter tuning, calibration,
routing, or sample-wise model selection. Final predictions and targets are
evaluated once in scaled space with the same masked global errors:
\begin{equation}
    \mathrm{MAE}=\frac{\sum_{b,j}m_{b,j}
    |\widetilde{\hat y}_{b,j}-\widetilde y_{b,j}|}
    {\sum_{b,j}m_{b,j}},
    \qquad
    \mathrm{MSE}=\frac{\sum_{b,j}m_{b,j}
    (\widetilde{\hat y}_{b,j}-\widetilde y_{b,j})^2}
    {\sum_{b,j}m_{b,j}}.
\end{equation}
Lower values are better. This protocol makes the comparison depend on the
forecasting model under the same observed-history and query masks, rather than
on different preprocessing, target definitions, or test-time adjustment.

\section{Main Experiments}
\label{sec:main-experiments}

\begin{table}[!htbp]
\caption{Standard-window benchmark comparison. Neural baselines report mean
$\pm$ standard deviation over three runs when available. \method{} reports mean
$\pm$ standard deviation over seeds 2024--2026. Lower is better.
\textbf{Bold} marks the best value and \underline{underlining} marks the
second-best value in each column.}
\label{tab:main-results}
\centering
\scriptsize
\resizebox{\textwidth}{!}{
\begin{tabular}{l|c|c|c|c}
\toprule
Method
& \multicolumn{2}{|c|}{HumanActivity}
& \multicolumn{2}{c}{USHCN} \\
\cmidrule(lr){2-3}
\cmidrule(lr){4-5}
& MSE & MAE & MSE & MAE \\
\midrule
PrimeNet & 4.2507$\pm$0.0041 & 1.7018$\pm$0.0011 & 0.4930$\pm$0.0015 & 0.4954$\pm$0.0018 \\
NeuralFlows & 0.1722$\pm$0.0090 & 0.3150$\pm$0.0094 & 0.2087$\pm$0.0258 & 0.3157$\pm$0.0187 \\
CRU & 0.1387$\pm$0.0073 & 0.2607$\pm$0.0092 & 0.2168$\pm$0.0162 & 0.3180$\pm$0.0248 \\
mTAN & 0.0993$\pm$0.0026 & 0.2219$\pm$0.0047 & 0.5561$\pm$0.2020 & 0.5015$\pm$0.0968 \\
SeFT & 1.3786$\pm$0.0024 & 0.9762$\pm$0.0007 & 0.3345$\pm$0.0022 & 0.4083$\pm$0.0084 \\
GNeuralFlow & 0.3936$\pm$0.1585 & 0.4541$\pm$0.0841 & 0.2205$\pm$0.0421 & 0.3286$\pm$0.0412 \\
GRU-D & 0.1893$\pm$0.0627 & 0.3253$\pm$0.0485 & 0.2097$\pm$0.0493 & 0.3045$\pm$0.0305 \\
Raindrop & 0.0916$\pm$0.0072 & 0.2114$\pm$0.0072 & 0.2035$\pm$0.0336 & 0.3029$\pm$0.0264 \\
Warpformer & 0.0446$\pm$0.0010 & 0.1219$\pm$0.0011 & 0.2058$\pm$0.0771 & 0.3089$\pm$0.0769 \\
tPatchGNN & 0.0438$\pm$0.0004 & 0.1228$\pm$0.0021 & \second{0.1596$\pm$0.0074} & \second{0.2668$\pm$0.0101} \\
GraFITi & 0.0439$\pm$0.0006 & 0.1226$\pm$0.0019 & 0.1669$\pm$0.0037 & 0.2730$\pm$0.0269 \\
APN & \best{0.0421$\pm$0.0001} & \second{0.1156$\pm$0.0002} & 0.1611$\pm$0.0205 & 0.2844$\pm$0.0520 \\
\midrule
\method{} & \second{0.0434$\pm$0.0008} & \best{0.1126$\pm$0.0008} & \best{0.1332$\pm$0.0083} & \best{0.2075$\pm$0.0067} \\
\bottomrule
\end{tabular}}

\vspace{0.75em}

\resizebox{\textwidth}{!}{
\begin{tabular}{l|c|c|c|c}
\toprule
Method
& \multicolumn{2}{|c|}{P12}
& \multicolumn{2}{c}{MIMIC} \\
\cmidrule(lr){2-3}
\cmidrule(lr){4-5}
& MSE & MAE & MSE & MAE \\
\midrule
PrimeNet & 0.7953$\pm$0.0000 & 0.6859$\pm$0.0001 & 0.9073$\pm$0.0001 & 0.6614$\pm$0.0001 \\
NeuralFlows & 0.4056$\pm$0.0033 & 0.4466$\pm$0.0027 & 0.6085$\pm$0.0101 & 0.5306$\pm$0.0066 \\
CRU & 0.6179$\pm$0.0045 & 0.5778$\pm$0.0031 & 0.5895$\pm$0.0092 & 0.5151$\pm$0.0048 \\
mTAN & 0.3809$\pm$0.0043 & 0.4291$\pm$0.0035 & 0.9408$\pm$0.1126 & 0.6755$\pm$0.0459 \\
SeFT & 0.7721$\pm$0.0021 & 0.6760$\pm$0.0029 & 0.9230$\pm$0.0015 & 0.6628$\pm$0.0008 \\
GNeuralFlow & 0.8207$\pm$0.0310 & 0.6759$\pm$0.0100 & 0.8957$\pm$0.0209 & 0.6450$\pm$0.0072 \\
GRU-D & 0.3419$\pm$0.0029 & 0.3992$\pm$0.0011 & 0.4759$\pm$0.0100 & 0.4526$\pm$0.0055 \\
Raindrop & 0.3478$\pm$0.0019 & 0.4044$\pm$0.0020 & 0.6754$\pm$0.1829 & 0.5444$\pm$0.0868 \\
Warpformer & \second{0.3059$\pm$0.0013} & 0.3655$\pm$0.0019 & \second{0.4302$\pm$0.0035} & \second{0.4025$\pm$0.0014} \\
tPatchGNN & 0.3117$\pm$0.0034 & 0.3676$\pm$0.0033 & 0.4431$\pm$0.0115 & 0.4077$\pm$0.0088 \\
GraFITi & 0.3070$\pm$0.0030 & \second{0.3633$\pm$0.0051} & 0.5035$\pm$0.0310 & 0.4442$\pm$0.0254 \\
APN & 0.3150$\pm$0.0016 & 0.3691$\pm$0.0032 & 0.4746$\pm$0.0045 & 0.4155$\pm$0.0063 \\
\midrule
\method{} & \best{0.3020$\pm$0.0014} & \best{0.3511$\pm$0.0005} & \best{0.4232$\pm$0.0022} & \best{0.3719$\pm$0.0043} \\
\bottomrule
\end{tabular}}
\end{table}

\begin{table}[!htbp]
\caption{\method{} standard-window seed stability. Entries are global scaled
errors for the three independent seeds used in Table~\ref{tab:main-results}.
Lower is better.}
\label{tab:adept-seeds}
\centering
\scriptsize
\resizebox{\textwidth}{!}{
\begin{tabular}{l|cc|cc|cc|cc}
\toprule
Seed
& \multicolumn{2}{c|}{HumanActivity}
& \multicolumn{2}{c|}{USHCN}
& \multicolumn{2}{c|}{P12}
& \multicolumn{2}{c}{MIMIC} \\
\cmidrule(lr){2-3}
\cmidrule(lr){4-5}
\cmidrule(lr){6-7}
\cmidrule(lr){8-9}
& MSE & MAE & MSE & MAE & MSE & MAE & MSE & MAE \\
\midrule
2024 & 0.0428 & 0.1134 & 0.1426 & 0.2132 & 0.3036 & 0.3514 & 0.4216 & 0.3710 \\
2025 & 0.0443 & 0.1118 & 0.1268 & 0.2093 & 0.3007 & 0.3514 & 0.4257 & 0.3765 \\
2026 & 0.0431 & 0.1127 & 0.1303 & 0.2001 & 0.3019 & 0.3505 & 0.4222 & 0.3680 \\
\midrule
Mean & 0.0434 & 0.1126 & 0.1332 & 0.2075 & 0.3020 & 0.3511 & 0.4232 & 0.3719 \\
Std. & 0.0008 & 0.0008 & 0.0083 & 0.0067 & 0.0014 & 0.0005 & 0.0022 & 0.0043 \\
\bottomrule
\end{tabular}}
\end{table}

\begin{table}[!htbp]
\caption{Entity-bootstrap uncertainty for \method{} at the standard benchmark
windows. Point estimates are target-weighted global scaled errors; 95\%
intervals resample seed--entity units with replacement.}
\label{tab:entity-bootstrap}
\centering
\scriptsize
\begin{tabular}{l|r|r|c|c}
\toprule
Dataset & Seed--entity units & Predictions & MAE [95\% CI] & MSE [95\% CI] \\
\midrule
P12 & 3,555 & 70,740 & 0.3511 [0.3456, 0.3568] & 0.3020 [0.2916, 0.3126] \\
MIMIC & 5,316 & 69,384 & 0.3719 [0.3648, 0.3793] & 0.4232 [0.3981, 0.4506] \\
USHCN & 336 & 1,104 & 0.2075 [0.1880, 0.2284] & 0.1332 [0.1093, 0.1618] \\
HumanActivity & 654 & 17,730 & 0.1126 [0.1073, 0.1183] & 0.0434 [0.0386, 0.0484] \\
\bottomrule
\end{tabular}
\end{table}

\begin{figure*}[!htbp]
\centering
\pgfplotsset{
    adeptWindowPanel/.style={
        width=0.445\linewidth,
        height=0.305\linewidth,
        xlabel={Lookback / Horizon},
        grid=both,
        grid style={gray!15},
        major grid style={gray!20},
        tick label style={font=\scriptsize},
        xticklabel style={rotate=30, anchor=east, yshift=-0.6ex},
        label style={font=\scriptsize},
        title style={font=\scriptsize, yshift=1.0ex},
        legend style={font=\scriptsize, draw=black, rounded corners=2pt, fill=none, at={(1.04,1.30)}, anchor=south, legend columns=4},
        scaled y ticks=false,
        clip=false
    },
    adeptModel/.style={black, thick, mark=*, mark size=1.8pt},
    apnStyle/.style={red!70!black, thick, mark=diamond*, mark size=2.5pt},
    grafitiStyle/.style={blue!70!black, thick, mark=triangle*, mark size=2.2pt},
    tpatchStyle/.style={green!50!black, thick, mark=square*, mark size=2pt}
}

\begin{subfigure}{\textwidth}
\centering
\begin{tikzpicture}
\begin{groupplot}[
    group style={group size=2 by 1, horizontal sep=1.80cm},
    adeptWindowPanel,
    xmin=0.7, xmax=4.3,
    xtick={1,2,3,4},
    xticklabels={{1000/100},{2000/200},{3000/300},{3000/600}}
]
\nextgroupplot[ylabel={MSE}, ymin=0.038, ymax=0.066]
\path[fill=gray!18, draw=none] (axis cs:2.65,0.038) rectangle (axis cs:3.35,0.066);
\addplot+[adeptModel] coordinates {(1,0.0565) (2,0.0494) (3,0.0436) (4,0.0505)};
\addlegendentry{\method{}}
\addplot+[apnStyle] coordinates {(1,0.0545) (2,0.0483) (3,0.0421) (4,0.0483)};
\addlegendentry{APN}
\addplot+[grafitiStyle] coordinates {(1,0.0603) (2,0.0514) (3,0.0439) (4,0.0506)};
\addlegendentry{GraFITi}
\addplot+[tpatchStyle] coordinates {(1,0.0620) (2,0.0531) (3,0.0470) (4,0.0524)};
\addlegendentry{tPatchGNN}
\nextgroupplot[ylabel={MAE}, ymin=0.103, ymax=0.145]
\path[fill=gray!18, draw=none] (axis cs:2.65,0.103) rectangle (axis cs:3.35,0.145);
\addplot+[adeptModel] coordinates {(1,0.1101) (2,0.1133) (3,0.1123) (4,0.1262)};
\addplot+[apnStyle] coordinates {(1,0.1139) (2,0.1159) (3,0.1156) (4,0.1265)};
\addplot+[grafitiStyle] coordinates {(1,0.1304) (2,0.1290) (3,0.1226) (4,0.1369)};
\addplot+[tpatchStyle] coordinates {(1,0.1322) (2,0.1311) (3,0.1323) (4,0.1398)};
\end{groupplot}
\end{tikzpicture}
\centerline{\footnotesize (a) HumanActivity}
\end{subfigure}

\vspace{0.95em}

\begin{subfigure}{\textwidth}
\centering
\begin{tikzpicture}
\begin{groupplot}[
    group style={group size=2 by 1, horizontal sep=1.80cm},
    adeptWindowPanel,
    xmin=0.7, xmax=4.3,
    xtick={1,2,3,4},
    xticklabels={{12/1},{24/3},{72/3},{72/6}}
]
\nextgroupplot[ylabel={MSE}, ymin=0.32, ymax=1.22]
\path[fill=gray!18, draw=none] (axis cs:2.65,0.32) rectangle (axis cs:3.35,1.22);
\addplot+[adeptModel] coordinates {(1,0.9514) (2,0.6764) (3,0.4207) (4,0.4261)};
\addplot+[apnStyle] coordinates {(1,0.9828) (2,0.7252) (3,0.4746) (4,0.4837)};
\addplot+[grafitiStyle] coordinates {(1,0.9915) (2,0.7319) (3,0.5035) (4,0.5357)};
\addplot+[tpatchStyle] coordinates {(1,1.1484) (2,0.7210) (3,0.4695) (4,0.5145)};
\nextgroupplot[ylabel={MAE}, ymin=0.33, ymax=0.64]
\path[fill=gray!18, draw=none] (axis cs:2.65,0.33) rectangle (axis cs:3.35,0.64);
\addplot+[adeptModel] coordinates {(1,0.5633) (2,0.5126) (3,0.3713) (4,0.3771)};
\addplot+[apnStyle] coordinates {(1,0.5982) (2,0.5558) (3,0.4155) (4,0.4278)};
\addplot+[grafitiStyle] coordinates {(1,0.6086) (2,0.5592) (3,0.4442) (4,0.4711)};
\addplot+[tpatchStyle] coordinates {(1,0.5791) (2,0.5517) (3,0.4080) (4,0.4311)};
\end{groupplot}
\end{tikzpicture}
\centerline{\footnotesize (b) MIMIC}
\end{subfigure}

\vspace{0.95em}

\begin{subfigure}{\textwidth}
\centering
\begin{tikzpicture}
\begin{groupplot}[
    group style={group size=2 by 1, horizontal sep=1.80cm},
    adeptWindowPanel,
    xmin=0.7, xmax=4.3,
    xtick={1,2,3,4},
    xticklabels={{12/1},{24/3},{36/3},{36/6}}
]
\nextgroupplot[ylabel={MSE}, ymin=0.265, ymax=0.345]
\path[fill=gray!18, draw=none] (axis cs:2.65,0.265) rectangle (axis cs:3.35,0.345);
\addplot+[adeptModel] coordinates {(1,0.3067) (2,0.3152) (3,0.3035) (4,0.3213)};
\addplot+[apnStyle] coordinates {(1,0.3249) (2,0.3213) (3,0.3150) (4,0.3345)};
\addplot+[grafitiStyle] coordinates {(1,0.3089) (2,0.3170) (3,0.3070) (4,0.3239)};
\addplot+[tpatchStyle] coordinates {(1,0.3143) (2,0.3229) (3,0.3117) (4,0.3280)};
\nextgroupplot[ylabel={MAE}, ymin=0.335, ymax=0.395]
\path[fill=gray!18, draw=none] (axis cs:2.65,0.335) rectangle (axis cs:3.35,0.395);
\addplot+[adeptModel] coordinates {(1,0.3476) (2,0.3569) (3,0.3525) (4,0.3698)};
\addplot+[apnStyle] coordinates {(1,0.3744) (2,0.3711) (3,0.3691) (4,0.3881)};
\addplot+[grafitiStyle] coordinates {(1,0.3568) (2,0.3706) (3,0.3633) (4,0.3813)};
\addplot+[tpatchStyle] coordinates {(1,0.3645) (2,0.3747) (3,0.3676) (4,0.3860)};
\end{groupplot}
\end{tikzpicture}
\centerline{\footnotesize (c) P12}
\end{subfigure}

\vspace{0.95em}

\begin{subfigure}{\textwidth}
\centering
\begin{tikzpicture}
\begin{groupplot}[
    group style={group size=2 by 1, horizontal sep=1.80cm},
    adeptWindowPanel,
    xmin=0.7, xmax=4.3,
    xtick={1,2,3,4},
    xticklabels={{50/1},{100/3},{150/3},{150/7}}
]
\nextgroupplot[ylabel={MSE}, ymin=0.12, ymax=0.78]
\path[fill=gray!18, draw=none] (axis cs:2.65,0.12) rectangle (axis cs:3.35,0.78);
\addplot+[adeptModel] coordinates {(1,0.6011) (2,0.5409) (3,0.1614) (4,0.4160)};
\addplot+[apnStyle] coordinates {(1,0.5989) (2,0.6457) (3,0.1611) (4,0.4103)};
\addplot+[grafitiStyle] coordinates {(1,0.5805) (2,0.6576) (3,0.1669) (4,0.4378)};
\addplot+[tpatchStyle] coordinates {(1,0.5504) (2,0.6109) (3,0.1596) (4,0.4067)};
\nextgroupplot[ylabel={MAE}, ymin=0.20, ymax=0.46]
\path[fill=gray!18, draw=none] (axis cs:2.65,0.20) rectangle (axis cs:3.35,0.46);
\addplot+[adeptModel] coordinates {(1,0.3517) (2,0.3211) (3,0.2239) (4,0.2765)};
\addplot+[apnStyle] coordinates {(1,0.4020) (2,0.4262) (3,0.2844) (4,0.3266)};
\addplot+[grafitiStyle] coordinates {(1,0.3976) (2,0.4394) (3,0.2730) (4,0.3539)};
\addplot+[tpatchStyle] coordinates {(1,0.4045) (2,0.4338) (3,0.2668) (4,0.3437)};
\end{groupplot}
\end{tikzpicture}
\centerline{\footnotesize (d) USHCN}
\end{subfigure}
\caption{Cross-model lookback--horizon comparison on a balanced four-window grid
per dataset. Each panel plots global scaled MSE and MAE for \method{} and the
three available neural baselines. Gray bands mark the standard benchmark
window. Lower is better.}
\label{fig:cross-model-window-sweep}
\end{figure*}

\paragraph{Standard-window comparison}
Table~\ref{tab:main-results} is the central comparison because all methods are
evaluated at the canonical benchmark windows. \method{} improves the strongest
standard-window MAE in every dataset. On USHCN, MAE falls from the best baseline
value of 0.2668 to 0.2075, and MSE falls from 0.1596 to 0.1332. This is the
largest relative improvement and is the clearest effect of the adaptive expert decoder:
history-derived seasonal and graph-aware summaries correct the climate regime
where the model without the adaptive expert decoder was window-sensitive. On MIMIC, \method{}
reduces MAE from 0.4025 to 0.3719 and MSE from 0.4302 to 0.4232. The MSE gain
is modest, but it is consistent across seeds and appears together with a larger
MAE improvement. On P12, \method{} improves MAE from 0.3633 to 0.3511 and MSE
from 0.3059 to 0.3020.

HumanActivity is the one standard-window case where \method{} does not win both
metrics. It obtains the best MAE, 0.1126 versus APN's 0.1156, but APN remains
lower on MSE, 0.0421 versus 0.0434. This should be read as a squared-error
trade-off, not as a complete loss on the dataset: the typical absolute error is
better, while a small number of larger deviations keep APN ahead in MSE. The
standard-window conclusion is therefore strong but bounded: \method{} leads all
four MAE blocks, leads three of four MSE blocks, and has one small but real MSE
counterexample.

\paragraph{Seed stability}
Table~\ref{tab:adept-seeds} shows that the standard-window result is not a
single lucky seed. P12 and HumanActivity are especially stable: the MAE standard
deviation is 0.0005 and 0.0008, respectively. MIMIC has slightly larger MAE
variation, 0.0043, but every seed remains well below the strongest baseline MAE.
USHCN is naturally less stable because it has far fewer masked targets, yet all
three \method{} seeds improve the baseline MAE by a large margin. The seed table also
clarifies why the paper does not claim HumanActivity MSE leadership: each seed
is close to APN, but the mean remains higher than APN's 0.0421.
Table~\ref{tab:entity-bootstrap} adds entity-level uncertainty. The intervals
are tightest for the larger clinical and activity benchmarks and widest for
USHCN, matching the smaller target count in the standard climate window.

\paragraph{Lookback--horizon robustness}
Figure~\ref{fig:cross-model-window-sweep} separates robust behavior from
single-window luck. \method{} wins all 16 MAE cells across the four datasets.
The strongest pattern is on MIMIC and P12, where \method{} wins every displayed
MSE and MAE cell. MIMIC gains are largest at the clinically useful longer
windows: at 72/3, MSE improves from the best baseline value of 0.4695 to 0.4207,
and MAE improves from 0.4080 to 0.3713. P12 is more compact but similarly
consistent, with small MSE margins and clear MAE margins at all four windows. USHCN changes substantially when the adaptive expert decoder is included. The full
\method{} model now wins all USHCN MAE cells and the 100/3 MSE cell, while
remaining slightly behind the best MSE at 50/1, 150/3, and 150/7. This is a
useful boundary: the adaptive experts greatly improve typical climate error,
but the best squared-error behavior at some climate horizons still belongs to a
baseline. HumanActivity has the opposite shape. \method{} wins every
HumanActivity MAE cell, but APN keeps the best MSE in all four cells. The sweep
therefore refines the headline claim: \method{} is uniformly strong on MAE,
dominant on P12 and MIMIC, greatly improved on USHCN, and still MSE-contested on
HumanActivity.

\section{Additional Experiments}
\label{sec:additional-experiments}

\begin{table}[!htbp]
\caption{Architecture diagnostics for \method{} at the standard benchmark
windows. Entries are means over seeds 2024--2026. ``No adaptive expert decoder''
removes the adaptive expert gate and history-derived scalar experts while
leaving the neural event graph decoder intact. Lower is better; \textbf{bold}
marks the best value and \underline{underlining} marks the second-best value
within each dataset--metric block.}
\label{tab:architecture-diagnostics}
\centering
{\tiny
\setlength{\tabcolsep}{3pt}
\renewcommand{\arraystretch}{0.92}
\resizebox{0.82\textwidth}{!}{
\begin{tabular}{l|l|c|c}
\toprule
Dataset & Variant & MSE & MAE \\
\midrule
\multirow{7}{*}{HumanActivity}
& \method{} & \best{0.0436} & \best{0.1124} \\
& No adaptive expert decoder & 0.0451 & 0.1232 \\
& No delta kernel & 0.0467 & \second{0.1229} \\
& No graph mixer & 0.0456 & 0.1239 \\
& No transformer & \second{0.0437} & \second{0.1229} \\
& No time encoding & 0.0451 & 0.1293 \\
& Direct decoder & 0.3098 & 0.4080 \\
\midrule
\multirow{7}{*}{USHCN}
& \method{} & 0.2335 & 0.2443 \\
& No adaptive expert decoder & \best{0.1410} & \best{0.2218} \\
& No delta kernel & 0.1540 & 0.2391 \\
& No graph mixer & 0.1627 & 0.2449 \\
& No transformer & \second{0.1434} & \second{0.2290} \\
& No time encoding & 0.1469 & 0.2593 \\
& Direct decoder & 0.2378 & 0.3451 \\
\midrule
\multirow{7}{*}{P12}
& \method{} & \second{0.3030} & \best{0.3509} \\
& No adaptive expert decoder & 0.3044 & \second{0.3616} \\
& No delta kernel & 0.3037 & 0.3625 \\
& No graph mixer & 0.3053 & 0.3666 \\
& No transformer & \best{0.3021} & \second{0.3616} \\
& No time encoding & 0.3079 & 0.3679 \\
& Direct decoder & 0.7936 & 0.6797 \\
\midrule
\multirow{7}{*}{MIMIC}
& \method{} & \best{0.4260} & \best{0.3734} \\
& No adaptive expert decoder & 0.4338 & 0.3904 \\
& No delta kernel & \second{0.4337} & \second{0.3805} \\
& No graph mixer & 0.4415 & 0.3903 \\
& No transformer & 0.4460 & 0.3916 \\
& No time encoding & 0.4417 & 0.3907 \\
& Direct decoder & 0.9632 & 0.6418 \\
\bottomrule
\end{tabular}}}
\end{table}

\begin{figure*}[!htbp]
\centering
\pgfplotsset{
    adeptTemporalPanel/.style={
        width=0.445\linewidth,
        height=0.305\linewidth,
        grid=both,
        grid style={gray!15},
        major grid style={gray!20},
        tick label style={font=\scriptsize},
        xticklabel style={rotate=35, anchor=east, yshift=-0.6ex},
        label style={font=\scriptsize},
        title style={font=\scriptsize, yshift=1.0ex},
        scaled y ticks=false,
        clip=false
    },
    adeptLine/.style={black, thick, mark=*, mark size=1.9pt}
}

\begin{subfigure}{\textwidth}
\centering
\begin{tikzpicture}
\begin{groupplot}[
    group style={group size=2 by 1, horizontal sep=1.80cm},
    adeptTemporalPanel,
    xtick={1,2,3,4,5,6,7,8,9},
    xticklabels={{Orig},{Shuffle},{Swap},{Drop},{Last},{Gaps},{75\%},{50\%},{25\%}}
]
\nextgroupplot[ylabel={MSE}, ymin=0.25, ymax=0.85]
\addplot+[adeptLine] coordinates {(1,0.3034) (2,0.4595) (3,0.5216) (4,0.3249) (5,0.3615) (6,0.7870) (7,0.3231) (8,0.3574) (9,0.4328)};
\nextgroupplot[ylabel={MAE}, ymin=0.32, ymax=0.72]
\addplot+[adeptLine] coordinates {(1,0.3509) (2,0.4580) (3,0.4910) (4,0.3641) (5,0.3798) (6,0.6799) (7,0.3658) (8,0.3892) (9,0.4348)};
\end{groupplot}
\end{tikzpicture}
\centerline{\footnotesize (a) P12}
\end{subfigure}

\vspace{0.95em}

\begin{subfigure}{\textwidth}
\centering
\begin{tikzpicture}
\begin{groupplot}[
    group style={group size=2 by 1, horizontal sep=1.80cm},
    adeptTemporalPanel,
    xtick={1,2,3,4,5,6,7,8,9},
    xticklabels={{Orig},{Shuffle},{Swap},{Drop},{Last},{Gaps},{75\%},{50\%},{25\%}}
]
\nextgroupplot[ylabel={MSE}, ymin=0.38, ymax=1.00]
\addplot+[adeptLine] coordinates {(1,0.4269) (2,0.5657) (3,0.6543) (4,0.4384) (5,0.4346) (6,0.9384) (7,0.4540) (8,0.5424) (9,0.6843)};
\nextgroupplot[ylabel={MAE}, ymin=0.34, ymax=0.66]
\addplot+[adeptLine] coordinates {(1,0.3713) (2,0.4519) (3,0.4923) (4,0.3778) (5,0.3784) (6,0.6399) (7,0.3920) (8,0.4374) (9,0.5106)};
\end{groupplot}
\end{tikzpicture}
\centerline{\footnotesize (b) MIMIC}
\end{subfigure}

\vspace{0.95em}

\begin{subfigure}{\textwidth}
\centering
\begin{tikzpicture}
\begin{groupplot}[
    group style={group size=2 by 1, horizontal sep=1.80cm},
    adeptTemporalPanel,
    xtick={1,2,3,4,5,6,7,8,9},
    xticklabels={{Orig},{Shuffle},{Swap},{Drop},{Last},{Gaps},{75\%},{50\%},{25\%}}
]
\nextgroupplot[ylabel={MSE}, ymin=0.12, ymax=0.38]
\addplot+[adeptLine] coordinates {(1,0.1467) (2,0.2346) (3,0.1902) (4,0.1473) (5,0.2622) (6,0.3359) (7,0.1570) (8,0.1532) (9,0.1696)};
\nextgroupplot[ylabel={MAE}, ymin=0.20, ymax=0.40]
\addplot+[adeptLine] coordinates {(1,0.2196) (2,0.2546) (3,0.2332) (4,0.2236) (5,0.3151) (6,0.3744) (7,0.2269) (8,0.2299) (9,0.2345)};
\end{groupplot}
\end{tikzpicture}
\centerline{\footnotesize (c) USHCN}
\end{subfigure}

\vspace{0.95em}

\begin{subfigure}{\textwidth}
\centering
\begin{tikzpicture}
\begin{groupplot}[
    group style={group size=2 by 1, horizontal sep=1.80cm},
    adeptTemporalPanel,
    xtick={1,2,3,4,5,6,7,8,9},
    xticklabels={{Orig},{Shuffle},{Swap},{Drop},{Last},{Gaps},{75\%},{50\%},{25\%}}
]
\nextgroupplot[ylabel={MSE}, ymin=0.03, ymax=4.5, ymode=log]
\addplot+[adeptLine] coordinates {(1,0.0433) (2,0.0859) (3,0.1027) (4,0.0437) (5,0.0548) (6,3.9569) (7,0.0476) (8,0.0506) (9,0.0719)};
\nextgroupplot[ylabel={MAE}, ymin=0.09, ymax=1.85, ymode=log]
\addplot+[adeptLine] coordinates {(1,0.1124) (2,0.1731) (3,0.1909) (4,0.1125) (5,0.1216) (6,1.6375) (7,0.1174) (8,0.1235) (9,0.1447)};
\end{groupplot}
\end{tikzpicture}
\centerline{\footnotesize (d) HumanActivity}
\end{subfigure}
\caption{Temporal perturbation and history-thinning tests for \method{} at the
standard benchmark windows, averaged over seeds 2024--2026. Orig uses unchanged
histories; Shuffle permutes timestamps; Swap exchanges early and late value
halves; Drop removes the first half of the history; Last keeps only the most
recent observation; Gaps keeps timestamps while replacing values by scaled zero;
and 75\%, 50\%, and 25\% retain random history subsets of the indicated size.
HumanActivity uses a log-scaled vertical axis because the time-gaps-only
intervention is much larger than the other interventions. Lower is better.}
\label{fig:adept-temporal-tests}
\end{figure*}

\begin{table*}[!htbp]
\caption{Cross-model temporal perturbation summary at the standard benchmark
windows. Deltas are relative to each model's original-history evaluation under
the same window and are aggregate, not paired, comparisons.}
\label{tab:cross-model-temporal}
\centering
\scriptsize
\resizebox{\textwidth}{!}{
\begin{tabular}{l|l|r|r|r|r|r|r}
\toprule
Dataset & Model & Orig. MSE & Shuffle $\Delta$MSE & Swap $\Delta$MSE & Gaps $\Delta$MSE & Orig. MAE & Gaps $\Delta$MAE \\
\midrule
\multirow{4}{*}{P12}
& \method{} & 0.3034 & +0.1561 & +0.2182 & +0.4835 & 0.3509 & +0.3290 \\
& APN & 0.3150 & +0.1410 & +0.1966 & +0.4816 & 0.3691 & +0.3168 \\
& GraFITi & 0.3070 & +0.1476 & +0.1905 & +0.4911 & 0.3633 & +0.3285 \\
& tPatchGNN & 0.3117 & +0.4565 & +0.1842 & +0.4843 & 0.3676 & +0.3178 \\
\midrule
\multirow{4}{*}{MIMIC}
& \method{} & 0.4269 & +0.1387 & +0.2274 & +0.5114 & 0.3713 & +0.2685 \\
& APN & 0.4746 & +0.1389 & +0.1890 & +0.4954 & 0.4155 & +0.2417 \\
& GraFITi & 0.5035 & +0.1126 & +0.1683 & +0.4672 & 0.4442 & +0.2294 \\
& tPatchGNN & 0.4695 & +0.1545 & +0.2108 & +0.4906 & 0.4080 & +0.2403 \\
\midrule
\multirow{4}{*}{USHCN}
& \method{} & 0.1467 & +0.0879 & +0.0435 & +0.1893 & 0.2196 & +0.1549 \\
& APN & 0.1611 & +0.0002 & +0.0024 & +0.0783 & 0.2844 & +0.0735 \\
& GraFITi & 0.1669 & +0.0007 & +0.0003 & +0.1032 & 0.2730 & +0.1035 \\
& tPatchGNN & 0.1596 & +0.0094 & +0.3894 & +0.1556 & 0.2668 & +0.1246 \\
\midrule
\multirow{4}{*}{HumanActivity}
& \method{} & 0.0433 & +0.0426 & +0.0595 & +3.9137 & 0.1124 & +1.5250 \\
& APN & 0.0421 & +0.0349 & +0.0581 & +4.0887 & 0.1156 & +1.5582 \\
& GraFITi & 0.0439 & +0.0345 & +0.0486 & +3.5535 & 0.1226 & +1.4151 \\
& tPatchGNN & 0.0470 & +0.3704 & +0.0527 & +3.7564 & 0.1323 & +1.4774 \\
\bottomrule
\end{tabular}}
\end{table*}

\paragraph{Architecture diagnostics}
Table~\ref{tab:architecture-diagnostics} separates the adaptive expert decoder
from the rest of the neural event graph. The no-adaptive-expert variant removes
only the expert gate and history-derived scalar experts, while the direct
decoder is a much more severe control that discards the main event-memory
pathway. The direct decoder is consistently poor, confirming that the model is
not just benefiting from dataset-level constants or simple target means. On
MIMIC, the full \method{} model is best on both metrics: MAE improves from
0.3904 to 0.3734 after adding the adaptive expert decoder, and MSE improves
from 0.4338 to 0.4260. On P12 and HumanActivity, the adaptive expert decoder
gives the best MAE, while no-transformer variants have slightly lower or
essentially tied MSE. This indicates that the expert decoder improves typical
absolute error, but the squared-error optimum can sometimes favor a less
expressive memory path. USHCN deserves a cautious reading. The standard \method{} benchmark in
Table~\ref{tab:main-results} is strong, but the architecture diagnostic run
shows that the no-adaptive-expert variant can be better under the diagnostic rerun. This
does not invalidate the standard table, which is the primary benchmark result;
it shows that USHCN is sensitive to small target counts and run context. The
important pattern is that learned event memory and graph-aware retrieval are
essential, while the adaptive gate improves USHCN strongly in the full benchmark
sweep but can overfit particular diagnostic splits.

\paragraph{Temporal perturbation and history thinning}
Figure~\ref{fig:adept-temporal-tests} tests whether \method{} actually uses
observed values and timing. The most severe intervention is replacing values by
scaled zero while preserving timestamps. This raises MSE from 0.3034 to 0.7870
on P12, from 0.4269 to 0.9384 on MIMIC, from 0.1467 to 0.3359 on USHCN, and
from 0.0433 to 3.9569 on HumanActivity. These degradations rule out a pure
time-mask shortcut.
Table~\ref{tab:cross-model-temporal} places the same temporal interventions
next to APN, GraFITi, and tPatchGNN. These are aggregate cross-model
comparisons rather than paired prediction-level tests, but they show that value
removal is strongly damaging for all models and that \method{} is especially
competitive on the original-history MAE columns.

Reordering also matters. Swapping early and late value halves increases MIMIC
MSE from 0.4269 to 0.6543 and HumanActivity MSE from 0.0433 to 0.1027.
Timestamp shuffling is especially harmful on P12 and HumanActivity, where the
model depends on correct temporal placement in addition to value identity.
History thinning is more gradual: keeping only 25\% of the history increases
MSE from 0.3034 to 0.4328 on P12, from 0.4269 to 0.6843 on MIMIC, from 0.1467
to 0.1696 on USHCN, and from 0.0433 to 0.0719 on HumanActivity. The smaller
USHCN thinning effect is consistent with the adaptive experts: once seasonal
and recent summaries are available, additional sparse climate points add less
than they do in the clinical datasets.

\paragraph{Additional-experiment synthesis}
The additional experiments support a specific mechanism. The transformer memory,
delta kernels, and graph retrieval provide the event representation; the
adaptive expert decoder then improves the final numerical forecast by deciding
when to trust same-variable persistence, graph-aware memory, local trend, or
the neural head.
This combination is especially useful for MAE: it wins every standard dataset
and every displayed window cell. The same diagnostics also show the boundary of
the approach. Squared-error leadership is still fragile on HumanActivity and
some USHCN windows, where small numbers of large deviations can outweigh better
typical absolute error. The model is therefore best described as broadly
state-of-the-art under the evaluated protocol, not universally dominant on every
metric and horizon.

\section{Discussion}
\label{sec:discussion}

\paragraph{Evidence}
The central modeling lesson is that IMTS forecasting should not separate
irregular representation from the final evidence-selection problem. A future
time--variable query often needs a local same-variable anchor, but this anchor
is not always sufficient: graph-neighbor measurements, entity-level state, and
short local trends become important when the requested variable is sparse or
when the horizon changes. \method{} addresses this by keeping observations as
events, retrieving them with query-conditioned attention and delta kernels, and
then using the adaptive expert decoder to decide how much numerical trust to
place in each historical source.

\paragraph{Robustness}
The empirical pattern is strongest for typical forecast accuracy. \method{}
achieves the best standard-window MAE on all four datasets and keeps the MAE
lead across every displayed lookback--horizon cell. This consistency is more
important than a single benchmark win: it suggests that the model is not merely
matched to one canonical window, but is repeatedly improving the ordinary
absolute error made on masked future targets. The MSE results are positive but
more conditional. \method{} leads standard-window MSE on P12, MIMIC, and
USHCN, yet APN remains better on HumanActivity MSE and several USHCN sweep
cells remain contested. Thus the main outcome is robust absolute-error
improvement, with squared-error sensitivity still visible when rare larger
deviations dominate the metric.

\paragraph{Mechanism}
The diagnostic experiments support the intended mechanism rather than a
shortcut explanation. Removing values while preserving timestamps sharply
degrades performance, so the model is not exploiting only time masks or target
availability patterns. Timestamp shuffling, early--late swaps, and history
thinning further show that temporal placement and history depth matter.
Ablations give the complementary architectural view: direct decoding is weak,
event memory and graph-aware retrieval are important, and the adaptive expert
decoder improves the final numerical forecast most reliably in MAE. Together,
these results indicate that the gains come from combining neural event memory
with adaptive, query-specific history anchors.

\paragraph{Protocol}
The comparison protocol also matters. \method{} does not use APN, GraFITi,
tPatchGNN, or any other baseline prediction as an expert. Its experts are
computed only from the observed history of the evaluated example, and their
gating weights are learned during training. Validation is used for checkpoint
and architecture selection, while test labels are reserved for final reporting.
This distinction is important because the adaptive decoder could otherwise be
mistaken for post hoc ensembling or test-time routing; in the implemented
model, it is an internal neural forecasting component.

\paragraph{Limitations}
The remaining weaknesses point to useful next steps. USHCN has few masked
targets at the standard window, so seed variation and diagnostic reversals are
expected. HumanActivity exposes a different issue: a model can improve MAE
while still losing MSE if occasional large deviations are not suppressed.
Future work should therefore study more outlier-aware residual control,
calibrated uncertainty for sparse query sets, and paired entity-level exports
for all baselines so that statistical comparisons can move beyond aggregate
tables. More broadly, the results suggest that the next progress in IMTS
forecasting may come less from choosing between graphs, patches, or continuous
dynamics, and more from designing decoders that know when each source of
historical evidence should be trusted.

\section{Conclusion}
\label{sec:conclusion}

\method{} frames irregular multivariate forecasting as a query-conditioned
event retrieval problem. The model keeps observations as time--variable events,
learns adaptive cross-variable structure, retrieves event memory with
continuous-time delta kernels, and decodes each future target through a gated
mixture of neural and history-derived experts. This design preserves strong
local forecasting anchors while remaining end-to-end neural and avoiding
external predictions, test labels, or post hoc routing. Across four IMTS benchmarks, \method{} delivers the best standard-window MAE on
every dataset and the best standard-window MSE on three of four, with a
remaining HumanActivity MSE gap that clarifies the method's squared-error
boundary. Window sweeps, perturbations, bootstrap intervals, and ablations show
that the improvement is not a mask shortcut or a persistence-only effect: it
depends on real values, correct timing, event-memory retrieval, graph context,
and adaptive expert selection. The resulting contribution is a
protocol-preserving forecaster whose main strength is robust absolute-error
accuracy across irregular clinical, climate, and activity time series.

\bibliographystyle{plainnat}
\bibliography{adept}

\clearpage
\appendix
\section{Notation}
\label{app:adept-symbols}

\begin{table}[!htbp]
\caption{Core notation used in the \method{} architecture description.}
\label{tab:adept-symbols}
\centering
\scriptsize
\setlength{\tabcolsep}{3pt}
\renewcommand{\arraystretch}{0.94}
\begin{tabular}{p{0.14\textwidth}p{0.32\textwidth}p{0.14\textwidth}p{0.32\textwidth}}
\toprule
Symbol & Description & Symbol & Description \\
\midrule
$b$ & Entity index. &
$\mathcal E_b$ & Irregular observed-event history for entity $b$. \\
$i,j,k,\ell$ & Event, query, delta-kernel, and expert indices. &
$\mathcal Q_b$ & Future time--variable query set for entity $b$. \\
$N_b,M_b$ & Numbers of observed events and future queries. &
$m_{b,j}$ & Active target-mask indicator for losses and metrics. \\
$t_i,\tau_j$ & Observed event time and future query time. &
$v_i,u_j$ & Observed variable and requested query variable. \\
$\widetilde x_i$ & Benchmark-scaled observed value. &
$\widetilde y_{b,j}$ & Benchmark-scaled target value. \\
\addlinespace
$h_i^{(0)}$ & Initial event token. &
$e_v,\bar e_v$ & Learned variable embedding and its normalized version. \\
$f_x,f_t$ & Value projection and Fourier time projection. &
$H_b$ & Transformer-contextualized event memory. \\
$d$ & Hidden width used in graph scaling. &
$A_{u,v}$ & Adaptive graph weight from query variable $u$ to source variable $v$. \\
$a_u$ & Projected graph context for query variable $u$. &
$K$ & Number of learned delta kernels. \\
$r_{j,k}$ & Delta-kernel memory context for query $j$ and kernel $k$. &
$\alpha_{j,i,k}$ & Normalized retrieval weight assigned to event $i$. \\
$s_k$ & Learned positive time scale for delta kernel $k$. &
$\beta,\gamma$ & Learned same-variable and graph-affinity bias strengths. \\
\addlinespace
$d_j$ & Deep query state before expert gating. &
$z^{\mathrm{nn}}_j$ & Neural base forecast. \\
$z^{\mathrm{last}}_j$ & Same-variable latest-value expert. &
$z^{\mathrm{ema}}_j$ & Same-variable exponential-memory expert. \\
$z^{\mathrm{trend}}_j$ & Local trend extrapolation expert. &
$z^{\mathrm{graph}}_j$ & Graph-aware cross-variable memory expert. \\
$z^{\mathrm{global}}_j$ & Global entity-memory expert. &
$\omega^u,\omega^g,\omega^0$ & Same-variable, graph-aware, and global expert weights. \\
$\widehat m_j$ & Clipped local slope estimate. &
$t^{\mathrm{last}}_{u_j}$ & Most recent observed time for query variable $u_j$. \\
$z_j,z_{j,\ell}$ & Full expert-value vector and its $\ell$th component. &
$\pi_j$ & Query-specific softmax gate over experts. \\
$W_\pi$ & Linear gate projection. &
$m_j$ & Same-variable history availability indicator. \\
$g_j$ & Query gap from the most recent observed event. &
$r_\theta$ & Bounded residual network. \\
$\widehat y_j$ & Final query forecast. &
$\mathcal L$ & Masked training objective. \\
\bottomrule
\end{tabular}
\end{table}

\section{ADEPT Configuration}
\label{app:adept-config}

\begin{table}[!htbp]
\caption{Global neural configuration of \method{} used in the benchmark runs.
Architecture diagnostics disable or replace one listed component at a time.}
\label{tab:adept-fixed-config}
\centering
\scriptsize
\begin{tabular}{p{0.24\textwidth}p{0.68\textwidth}}
\toprule
Component & Fixed setting \\
\midrule
Event representation & Each finite historical observation is an event token with
scaled value projection, learned variable embedding, and Fourier time features. \\
\addlinespace
Event memory & Transformer encoder with batch-first self-attention over padded
irregular event sets; default hidden width 96, two encoder layers, four
attention heads, and dropout 0.08--0.10 depending on run family. \\
\addlinespace
Adaptive variable graph & Graph weights are learned from normalized variable
embeddings through a softmax similarity matrix and are used in query
construction, graph-aware expert summaries, and temporal retrieval. \\
\addlinespace
Query decoder & Each future target is represented as a time--variable query.
The query attends to event memory with multi-head cross-attention and receives
graph context and normalized time-gap features. \\
\addlinespace
Delta kernels & Six learned positive time scales retrieve event memory using
recency, same-variable bias, and adaptive graph strength in the \method{} benchmark
runs. The concatenated kernel contexts are projected before decoding. \\
\addlinespace
Adaptive experts & The final head gates between the neural forecast, same-variable
latest value, same-variable exponential memory, local trend extrapolation,
graph-aware cross-variable memory, and global entity memory. Expert
values, a same-variable availability indicator, and the query gap are passed to
the gate. \\
\addlinespace
Loss and optimization & End-to-end masked loss combining MSE, smooth-$L_1$, and
MAE terms with weights 0.50, 0.20, and 0.30; AdamW with learning rate
$6\times10^{-4}$ to $8\times10^{-4}$ and weight decay $10^{-4}$; cosine
learning-rate decay. \\
\addlinespace
Validation protocol & Validation split is used for early stopping with patience
12--16. The \method{} validation score is MSE plus 0.25 times MAE. No test labels are
used for routing, calibration, or sample-wise model selection. \\
\addlinespace
Architectural diagnostics & The diagnostic suite compares the full \method{}
model with no-adaptive-expert-decoder, no-delta-kernel, no-graph-mixer,
no-transformer, no-time-encoding, and direct-decoder ablations, all trained and
evaluated through the same runner. \\
\addlinespace
Event limits & Long histories are capped by retaining recent events and, during
training, a deterministic random subset of older events. \method{} benchmark runs use
up to 2048 events per entity window and cap query batches for memory control. \\
\bottomrule
\end{tabular}
\end{table}

\end{document}
